\documentclass[12pt]{ximera}
\title{Homework 7: Counting --- Multiplication Principle, Permutations, Combinations}
\begin{document}
\begin{abstract}
Sections 8.1--8.2: deciding which counting technique a situation calls for --- the
multiplication principle, permutations, or combinations --- and then carrying out the
count.
\end{abstract}
\maketitle

\section*{Sections 8.1--8.2: Counting Techniques}

The question to ask in every one of these problems is: \emph{does order matter, and
may things repeat?}
\begin{itemize}
\item Order matters, no repeats, choosing $r$ from $n$: permutations,
      $P(n,r)=\frac{n!}{(n-r)!}$.
\item Order does not matter, no repeats: combinations,
      $C(n,r)=\frac{n!}{r!\,(n-r)!}$.
\item Independent choices made in sequence, repeats allowed: multiplication
      principle.
\end{itemize}

\begin{problem}
For each of the following situations, use the multiplication principle, the
permutations formula, or the combinations formula to count appropriately.

\begin{prompt}
\textbf{(a)} How many ways are there for a track coach to select athletes for the
first, second, third, and fourth legs of a $4\times100\text{m}$ relay from a team of
$12$ runners?

Which technique applies?
\begin{multipleChoice}
\choice[correct]{Permutations --- order matters and no runner repeats}
\choice{Combinations --- order does not matter}
\choice{Multiplication principle with repetition allowed}
\end{multipleChoice}

The count is $\answer{11880}$.

\begin{feedback}
The four legs are distinguishable, so the order of selection matters, and no runner
can run two legs. This is a permutation of $4$ runners from $12$:
\[
P(12,4)=\frac{12!}{8!}=12\cdot 11\cdot 10\cdot 9 = 11{,}880.
\]
\end{feedback}
\end{prompt}

\begin{prompt}
\textbf{(b)} How many different ways can $4$ professors be selected to teach Finite
Mathematics from among $10$ teaching professors?

Which technique applies?
\begin{multipleChoice}
\choice{Permutations --- order matters}
\choice[correct]{Combinations --- we only care which four are chosen}
\choice{Multiplication principle with repetition allowed}
\end{multipleChoice}

The count is $\answer{210}$.

\begin{feedback}
Here we only care \emph{which} four professors are chosen, not in what order, and
nobody is chosen twice. That is a combination:
\[
C(10,4)=\frac{10!}{4!\,6!}=\frac{10\cdot 9\cdot 8\cdot 7}{4\cdot 3\cdot 2\cdot 1}=210.
\]
\end{feedback}
\end{prompt}

\begin{prompt}
\textbf{(c)} How many ways can each of the $7$ sections of Finite Mathematics be
assigned one of $4$ particular teaching professors (allowing professors to teach
multiple sections)?

Which technique applies?
\begin{multipleChoice}
\choice{Permutations --- order matters and nothing repeats}
\choice{Combinations --- order does not matter}
\choice[correct]{Multiplication principle --- each section is an independent choice, repeats allowed}
\end{multipleChoice}

The count is $\answer{16384}$.

\begin{hint}
Go section by section. How many professors could be assigned to section 1? To
section 2, given that professors may repeat?
\end{hint}

\begin{feedback}
Each of the $7$ sections independently gets any one of the $4$ professors, and
professors may be reused. By the multiplication principle,
\[
4\cdot 4\cdot 4\cdot 4\cdot 4\cdot 4\cdot 4 = 4^7 = 16{,}384.
\]
Note the base and exponent: it is $4^7$, not $7^4$. Each of the $7$ \emph{sections}
is a decision, and each decision has $4$ options.
\end{feedback}
\end{prompt}

\begin{prompt}
\textbf{(d)} How many ways are there to guess the birth order of the ten siblings in
my grampa's family growing up?

The count is $\answer{3628800}$.

\begin{feedback}
This is a full ordering of all $10$ siblings, which is $10$ factorial:
\[
P(10,10)=10! = 3{,}628{,}800.
\]
\end{feedback}
\end{prompt}
\end{problem}

\begin{problem}
Use the multiplication principle and/or permutations to count the following.

Suppose you want to plan outfits for a $5$-day trip. You've made a checklist of the
decisions you need to make:
\begin{itemize}
    \item You own $11$ tops, and need to select and order $5$ for the trip.
    \item You own $7$ bottoms, and similarly need to select and order $5$ for the trip.
    \item You own $9$ pairs of socks, and need to select and order $5$ pairs for the trip.
\end{itemize}
How many possible outfit arrangements are there for the week?

\begin{prompt}
Write out the formula you would use. (You do not need to calculate it fully.)

\begin{freeResponse}
\end{freeResponse}

\begin{hint}
Each of the three checklist items is a permutation --- you are selecting \emph{and
ordering}, since assigning a top to Monday versus Friday gives a different week.
Then combine the three independent decisions.
\end{hint}

\begin{feedback}
Each line of the checklist is a permutation, because the items are assigned to
particular days:
\[
P(11,5)\cdot P(7,5)\cdot P(9,5)
= \frac{11!}{6!}\cdot\frac{7!}{2!}\cdot\frac{9!}{4!}.
\]
The three decisions are independent, so by the multiplication principle they
multiply. Numerically this is
\[
55{,}440 \cdot 2{,}520 \cdot 15{,}120 = 2{,}112{,}397{,}056{,}000,
\]
a little over two trillion possible weeks.
\end{feedback}
\end{prompt}

\begin{prompt}
To check the pieces, enter each permutation's value:

$P(11,5)=\answer{55440}$ \qquad
$P(7,5)=\answer{2520}$ \qquad
$P(9,5)=\answer{15120}$

\begin{feedback}
\[
P(11,5)=11\cdot 10\cdot 9\cdot 8\cdot 7 = 55{,}440,
\]
\[
P(7,5)=7\cdot 6\cdot 5\cdot 4\cdot 3 = 2{,}520,
\]
\[
P(9,5)=9\cdot 8\cdot 7\cdot 6\cdot 5 = 15{,}120.
\]
\end{feedback}
\end{prompt}
\end{problem}

\begin{problem}
Use our counting techniques to answer the following.

In the 2024 Olympics, Team USA was represented in the team finals for women's
gymnastics by $5$ of their athletes. Women's gymnastics has $4$ different apparatuses
(vault, uneven bars, balance beam, and floor), and for each apparatus, the coach
selects $3$ gymnasts to represent the team on that apparatus.

\begin{prompt}
\textbf{(a)} How many ways were there to select which of these gymnasts would
represent Team USA on the vault?

$\answer{10}$

\begin{feedback}
Choosing $3$ gymnasts out of $5$, with order irrelevant:
\[
C(5,3)=\frac{5!}{3!\,2!}=\frac{5\cdot 4}{2\cdot 1}=10.
\]
\end{feedback}
\end{prompt}

\begin{prompt}
\textbf{(b)} How many ways were there to select representatives for all four
apparatuses, altogether?

$\answer{10000}$

\begin{hint}
The same gymnast may compete on more than one apparatus, so the four selections do
not interfere with each other. Each apparatus is an independent repeat of part (a).
\end{hint}

\begin{feedback}
Each apparatus independently gets one of the $C(5,3)=10$ possible trios, and a
gymnast may appear on several apparatuses. By the multiplication principle,
\[
10 \cdot 10 \cdot 10 \cdot 10 = 10^4 = 10{,}000.
\]
\end{feedback}
\end{prompt}
\end{problem}

\end{document}
